% Part: proof-theory
% Chapter: proof-search
% Section: search-algorithm
%READERNOTE{SOL6-A312: सर्व योग्य पदे वापरून झाल्यावर संबंधित अरिक्त संचातील स्थिरांक पुन्हा वापरायचा; संचाबाहेरचा स्थिरांक आणून पुढील आयगेन-चर अट मोडायची नाही.}

\documentclass[../../../include/open-logic-section]{subfiles}

\begin{document}

\olfileid{pt}{ps}{sal}

\olsection{\Log{G3c} साठी सिद्धता-शोध अल्गोरिदम}

\Log{G3c} साठी सिद्धता-शोध अल्गोरिदम मांडू.
हा एकमेव शक्य अल्गोरिदम नाही आणि सर्वांत कार्यक्षमही
नाही. पण तो बरोबर आहे:~$\Gamma \Sequent \Delta$
या क्रमवर्तीची !!a{proof} असेल, तर अखेरीस तो
थांबतो आणि ती !!a{proof} शोधतो. तो विश्वासार्हही
आहे: !!a{proof} न सापडता तो थांबला, किंवा कधीच
समाप्त झाला नाही, तर~$\Gamma \Sequent \Delta$
मुळातच वैध नाही.
या बंद-पद शोधात सुरुवातीच्या क्रमवर्तीतील सर्व सूत्रे
वाक्ये, म्हणजे मुक्त चल नसलेली सूत्रे, असावीत.

अल्गोरिदमचा अत्यावश्यक गुणधर्म असा की
$\Gamma \Sequent \Delta$ मधील किंवा सिद्धता-शोधात
निर्माण झालेल्या क्रमवर्तीतील प्रत्येक !!{formula}
न्यूनीकृत होते: उलट्या दिशेने लावलेल्या
अनुमाननियमात त्याला मुख्य सूत्र मानले जाते,
आणि गरजेइतक्या वेळा असे केले जाते.
या गुणधर्माला \emph{न्याय्यता} म्हणू.

अल्गोरिदम टप्प्याटप्प्याने चालतो.
प्रत्येक टप्प्याचे फलित क्रमवर्तींचा वृक्ष असतो.
पुढील टप्प्याचे फलित मिळवताना आधीच स्वयंसिद्धक
नसलेल्या प्रत्येक सर्वांत वरच्या क्रमवर्तीतील
!!a{formula} चे न्यूनीकरण करतो.

~$0$ व्या टप्प्यावर~$\Gamma \Sequent \Delta$
ने सुरुवात करतो. न्याय्यता सुनिश्चित करण्यासाठी
$\Gamma$ आणि~$\Delta$ मधील !!{formula}s च्या
प्रत्येक आस्थितीला, म्हणजे वेगवेगळ्या !!{formula}s
च्या आस्थितींना, वेगळा संख्यात्मक निर्देशांक
देतो. उदाहरणार्थ, !!{formula}s डावीकडून उजवीकडे
मोजा. या संख्या ऊर्ध्वनिर्देशांक म्हणून लिहितो.
उदाहरणार्थ,~$!A, !B \Sequent !C$ ची !!a{proof}
शोधायची असेल, तर~$0$ व्या टप्प्याचे फलित
$!A^1, !B^2 \Sequent !C^3$ असते. अशा निर्देशांकित
क्रमवर्तीतील ऊर्ध्वनिर्देशांकांपैकी सर्वांत मोठा
~$n$ हा तिचा \emph{कमाल निर्देशांक}
(उदाहरणात~$3$) आहे.

क्रमवर्तीत संख्यापक असतील, तर
उजवीकडील~$\lexists[x][!A(x)]$ किंवा डावीकडील
$\lforall[x][!A(x)]$ न्यूनीकृत करताना कोणती पदे
वापरायची तेही ठरवावे लागते. ~$\Gamma$ आणि~$\Delta$
मधील फलनचिन्हे वापरून~!!{constant}s
$\Obj c_0$, $\Obj c_1$, $\Obj c_2$, \dots
यांपासून बनणारी पदेच लागतात. या !!{constant}s
पासून बनणाऱ्या सर्व पदांचा
एक निश्चित क्रम दिला आहे असे मानू; म्हणजे,
उदाहरणार्थ,~``$\Gamma \Sequent \Delta$ मध्ये
न आलेला पहिला !!{constant}'' असे म्हणता येईल.
अल्गोरिदमने बनवलेल्या वृक्षातील प्रत्येक
निर्देशांकित क्रमवर्तीशी !!{constant}s चा एक संच
संबंधित असतो. ~$0$ व्या टप्प्यातील क्रमवर्तीशी
तिच्यात येणाऱ्या सर्व !!{constant}s चा संच
(असतील तर) संबंधित असतो; त्यात एकही
!!{constant}s नसतील, तर~$\{\Obj c_0\}$ असतो.

~$n$ व्या टप्प्यावर संबंधित !!{constant}s
च्या संचांसह निर्देशांकित क्रमवर्तींचा वृक्ष
अल्गोरिदमने निर्माण केला आहे असे मानू.
~$n+1$ व्या टप्प्यावर प्रत्येक सर्वांत वरच्या
क्रमवर्तीसाठी पुढीलपैकी योग्य ती कृती करतो.

~$!A$ हे !!a{formula} क्रमवर्तीच्या डाव्या
आणि उजव्या दोन्ही बाजूंना असेल, किंवा डावीकडे
$\lfalse$ असेल, तर काही करत नाही: अशा क्रमवर्तींच्या
~\Log{G3c} मधील !!{proof}s असतात; या विशिष्ट
सर्वांत वरच्या क्रमवर्तीसाठी पुढे काही करायचे नसते.

क्रमवर्तीत अणुरूप नसलेले एकही !!{formula}
नसेल, तर अल्गोरिदम अयशस्वी म्हणून थांबवा.
$\Gamma \Sequent \Delta$ ची सिद्धता नाही.

अन्यथा सर्वांत लहान निर्देशांक~$i$ असलेले
अणुरूप नसलेले !!{formula}~$!A$ निवडा. मग विचाराधीन
सर्वांत वरची क्रमवर्ती~$!A^i, \Pi \Sequent \Lambda$
किंवा~$\Pi \Sequent \Lambda, !A^i$ अशी असते.
त्या क्रमवर्तीचा कमाल निर्देशांक~$k$ आणि
तिच्याशी संबंधित !!{constant}s चा संच~$C$ असू द्या.

~$!A^i$ चे ``न्यूनीकरण'' करू:~$!A$
चा !!{main operator} आणि~$!A^i$ डावीकडे आहे
की उजवीकडे, यांवर ठरणाऱ्या अनुमाननियमानुसार
नव्या सर्वांत वरच्या क्रमवर्ती तयार करतो.

\begin{enumerate}
  \item $!A \ident !B \land !C$ असून ते डावीकडे
  असेल, तर~$(!B \land !C)^i, \Pi \Sequent \Lambda$
  च्या वर एक नवी सर्वोच्च क्रमवर्ती जोडा:
  \[
  \Axiom$!B^{k+1}, !C^{k+2}, \Pi \fCenter \Lambda$
  \RightLabel{\LeftR{\land}}
  \UnaryInf$(!B \land !C)^i, \Pi \fCenter \Lambda$
  \DisplayProof
  \]
  नव्या क्रमवर्तीशी संबंधित !!{constant}s चा संच~$C$ आहे.
  \item $!A \ident !B \land !C$ असून ते उजवीकडे
  असेल, तर~$\Pi \Sequent \Lambda, (!B \land !C)^i$
  च्या वर दोन नव्या सर्वोच्च क्रमवर्ती जोडा:
  \[
  \Axiom$\Pi \fCenter \Lambda, !B^{k+1}$
  \Axiom$\Pi \fCenter \Lambda, !C^{k+1}$
  \RightLabel{\RightR{\land}}
  \BinaryInf$\Pi \fCenter \Lambda, (!B \land !C)^i$
  \DisplayProof
  \]
  प्रत्येक नव्या क्रमवर्तीशी संबंधित !!{constant}s चा संच~$C$ आहे.

  \item $!A \ident \lnot !B$,
  $!A \ident !B \lor !C$ आणि~$!A \ident !B \lif !C$
  हे प्रसंगही याच प्रकारे हाताळता येतात;
  ते सरावासाठी ठेवले आहेत.

  \item $!A \ident \lexists[x][!B(x)]$ असून ते
  डावीकडे असेल, तर~$\lexists[x][!B(x)]^i,
  \Pi \Sequent \Lambda$ च्या वर एक नवी
  सर्वोच्च क्रमवर्ती जोडा:
  \[
  \Axiom$!B(c)^{k+1}, \Pi \fCenter \Lambda$
  \RightLabel{\LeftR{\lexists}}
  \UnaryInf$\lexists[x][!B(x)]^i, \Pi \fCenter \Lambda$
  \DisplayProof
  \]
  येथे~$c$ हा~$C$ मध्ये नसलेला पहिला !!{constant}
  आहे. नव्या क्रमवर्तीशी संबंधित !!{constant}s
  चा संच~$C \cup \{c\}$ आहे.
  \item $!A \ident \lexists[x][!B(x)]$ असून ते
  उजवीकडे असेल, तर~$\Pi \Sequent \Lambda,
  \lexists[x][!B(x)]^i$ च्या वर एक नवी
  सर्वोच्च क्रमवर्ती जोडा:
  \[
  \Axiom$\Pi \fCenter \Lambda, \lexists[x][!B(x)]^{k+1}, !B(t)^{k+2}$
  \RightLabel{\RightR{\lexists}}
  \UnaryInf$\Pi \fCenter \Lambda, \lexists[x][!B(x)]^i$
  \DisplayProof
  \]
  येथे~$t$ हे~$C$ मधीलच !!{constant}s असलेले,
  आणि या क्रमवर्तीखाली उजवीकडील
  $\lexists[x][!B(x)]$ च्या आधीच्या न्यूनीकरणात
  न वापरलेले पहिले पद आहे. अशी सर्व पदे वापरून
  झाली असतील, तर~$C$ मधील पहिला स्थिरांक पुन्हा वापरा.
  $C$ सुरुवातीपासून अरिक्त आहे आणि पुढे त्यात फक्त भर
  पडते; म्हणून हा स्थिरांक उपलब्ध आहे आणि~$C$ बाहेरचा
  नवा स्थिरांक नकळत क्रमवर्तीत येत नाही.
  नव्या क्रमवर्तीशी संबंधित !!{constant}s चा संच~$C$.
  \item $!A \ident \lforall[x][!B(x)]$
  हे प्रसंगही याच प्रकारे हाताळता येतात;
  ते सरावासाठी ठेवले आहेत.
\end{enumerate}

~$n+1$ व्या टप्प्यावर कोणत्याही सर्वांत वरच्या
क्रमवर्तीच्या वर नवी क्रमवर्ती जोडली गेली नसेल,
तर अल्गोरिदम यशस्वीपणे संपतो. अशा वेळी
~$n+1$ व्या टप्प्यातील वृक्षाचे सर्व निर्देशांक
काढल्यावर~$\Gamma \Sequent \Delta$ ची
!!a{proof} मिळते. सर्वांत वरच्या क्रमवर्तींचे
रूप~$!A, \Pi \Sequent \Lambda, !A$ किंवा
$\lfalse, \Pi \Sequent \Lambda$ असे असते.
सर्व अनुमाने~\Log{G3c} मधील योग्य अनुमाने असतात.
विधानीय अनुमानांसाठी हे स्पष्ट आहे. प्रत्येक
टप्प्यावर क्रमवर्तीशी संबंधित !!{constant}s च्या
संच~$C$ मध्ये त्या क्रमवर्तीतील सर्व !!{constant}s असतात.
म्हणून~\LeftR{\lexists} आणि~\RightR{\lforall}
यांनी न्यूनीकरण करताना आयगेन-चराची अट पूर्ण होते:
आयगेन-चर~$c$ हा निष्कर्षाशी संबंधित संच~$C$
मध्ये नसलेला !!a{constant} असल्याने तो
~$c$ निष्कर्षातही आलेला नव्हता. म्हणून:

\begin{prop}
~\Log{G3c} साठीचा सिद्धता-शोध अल्गोरिदम बरोबर आहे:
तो यशस्वीपणे संपला, तर सुरुवातीच्या
क्रमवर्ती~$\Gamma \Sequent \Delta$ ची
!!a{proof} असते.
\end{prop}

\end{document}
